\documentclass[11pt]{article}
\usepackage[T1]{fontenc}
\usepackage[utf8]{inputenc}
\usepackage[margin=27mm]{geometry}
\usepackage{amsmath,amssymb,amsthm}
\usepackage[hidelinks]{hyperref}
\newtheorem{theorem}{Theorem}
\newtheorem{lemma}[theorem]{Lemma}
\newtheorem{corollary}[theorem]{Corollary}
\newcommand{\tr}{\operatorname{tr}}
\newcommand{\rank}{\operatorname{rank}}
\newcommand{\op}{\mathrm{op}}
\newcommand{\F}{\mathrm{F}}
\title{An exponential uniqueness bound for Gaussian SUR with equation-specific matrix designs}
\author{OpenAI gpt-6.1-sol, reasoning configuration ultra}
\date{28 October 2026}
\begin{document}
\maketitle

\begin{abstract}
Consider a correctly specified Gaussian seemingly unrelated regression
system with unrestricted positive-definite error covariance and an
arbitrary full-column-rank design in each equation. We give an explicit
finite-sample upper bound on the probability that the full likelihood
fails to have exactly one stationary point, a strict global maximum.
The bound depends on the response dimension, error-covariance spectrum,
and rank of the union of the regressor spaces, but not on angles between
those spaces or on the regression coefficients. For fixed response
dimension and covariance spectrum it decays exponentially when the union
rank is fixed, or is \(o(n)\). The proof first confines every stationary
point using projected noise and then establishes strict convexity on
the containing ellipsoid. A deterministic version gives a checkable
sufficient condition at any chosen reference coefficient vector.
\end{abstract}

\section{Original target and checked prior scope}

Question 11 of the AIM workshop list pairs classical MANOVA unimodality
with Takemura's question about the speed at which bivariate SUR
likelihood becomes unimodal \cite{aim}. The current repository statement
\cite{statement} preserves that question without specifying a sampling
law or a quantitative event. Its existing attempt \cite{attempt}
proves common-design Gaussian MANOVA uniqueness, but supplies no SUR
sampling rate.

Drton and Richardson \cite{dr} show that a simple bivariate SUR likelihood
can have multiple modes and that their probability vanishes
asymptotically. A September 2026 unrefereed preprint by Ferudun
\cite{ferudun} already gives explicit exponential bounds for the crossed
bivariate model with one private regressor in each equation, also
allowing common regressors. Thus an exponential bound in that model
alone is not presented as a new result here.

The contribution proved below covers arbitrary equation-specific
\emph{matrix} designs and any finite response dimension. It includes
several private regressors per equation, which the stated scope of
\cite{ferudun} excludes. Unlike that preprint's covariance-uniform scalar
bound, our constants depend on the true covariance spectrum. Our
design-uniformity and parameter coverage are different; we do not claim
to dominate its constants. No claim of exhaustive literature-wide
priority is made. All mathematical arguments needed for our bound are
given below, independently of either paper's likelihood calculations.

\section{Model and finite-sample statement}

For \(j=1,\ldots,q\), let
\[
 Y_j=X_j\beta_j^0+E_j,\qquad
 X_j\in\mathbb R^{n\times p_j},\quad
 \beta_j^0\in\mathbb R^{p_j}.
\]
Each \(X_j\) is deterministic and has full column rank. Zero-column
designs are allowed, with the corresponding coefficient space trivial.
Assume the total coefficient dimension is positive. The rows of
\(E=[E_1,\ldots,E_q]\) are independent \(N_q(0,\Sigma_0)\), where
\(\Sigma_0\) is positive definite. Both the coefficients and the
unrestricted covariance \(\Sigma\succ0\) are likelihood parameters.
There is no covariance sparsity constraint or missingness.

Let \(P\) be the orthogonal projection onto
\(\mathcal U=\sum_j\operatorname{col}(X_j)\), and put \(r=\dim\mathcal U\).
Write
\[
 \lambda_0=\lambda_{\min}(\Sigma_0),\quad
 \Lambda_0=\lambda_{\max}(\Sigma_0),\quad
 \lambda=\lambda_0/2,\quad \Lambda=3\Lambda_0/2,
 \qquad
 b_*=\frac{\lambda^3}{1024q(\Lambda+\lambda)^2}.
 \tag{1}
\]
All covariance matrices below use normalization \(1/n\).
A stationary point refers to derivatives in all coefficient coordinates
and all independent symmetric covariance coordinates on the open
parameter domain.

\begin{theorem}[Explicit Gaussian SUR uniqueness bound]\label{thm:main}
Under these assumptions the probability that the full likelihood does
not have exactly one stationary point, which is a strict global
maximum, is at most
\[
 \min\left\{1,\;
 2\,9^q e^{-n/128}
 +2^{rq/2}\exp\left(-\frac{n b_*}{4\Lambda_0}\right)
 \right\}.
 \tag{2}
\]
The bound is uniform over all true coefficient vectors and all
equation-specific full-column-rank designs with the specified union
rank. For fixed \(q,\lambda_0,\Lambda_0\) and \(r_n=o(n)\), its
right-hand side is at most \(C e^{-cn}\) for all sufficiently large
\(n\), for positive constants \(c,C\). The same holds in particular
for fixed union rank.
\end{theorem}

This is a bound for the entire unbounded coefficient domain and the
entire positive-definite covariance domain. It does not describe only
the stationary points reached from a particular initialization.
For small \(n\), (2) can be vacuous. A finite-dimensional Gaussian
model and the displayed probability event provide a precise answer
to one sampling interpretation of the original SUR question.

\section{A deterministic global uniqueness certificate}

The deterministic argument works at an arbitrary reference coefficient
vector. Let \(E\) denote its residual matrix; it need not be random
in this section. For coefficient displacement
\(\delta=(\delta_1,\ldots,\delta_q)\), set
\[
 A(\delta)=[X_1\delta_1,\ldots,X_q\delta_q],\quad
 D(\delta)=A(\delta)^{\mathsf T}A(\delta)/n,\quad
 C(\delta)=E^{\mathsf T}A(\delta)/n,
\]
\[
 \Omega=E^{\mathsf T}E/n,\qquad B=E^{\mathsf T}PE/n,\qquad
 S(\delta)=\Omega-C(\delta)-C(\delta)^{\mathsf T}+D(\delta).
 \tag{3}
\]
Thus \(S(\delta)\) is the residual covariance at the displaced
coefficient vector. Write \(t(\delta)=\tr D(\delta)\).
Full column rank makes \(t\) a positive-definite quadratic form on
the total coefficient space.

\begin{lemma}[Projected-noise certificate]\label{lem:certificate}
Suppose \(0<\lambda\leq\Lambda\),
\[
 \lambda I_q\preceq\Omega\preceq\Lambda I_q,\qquad
 0\preceq B\preceq b I_q,\qquad
 b\leq\frac{\lambda^3}{1024q(\Lambda+\lambda)^2}.
 \tag{4}
\]
Then \(S(\delta)\succ0\) for every \(\delta\), and
\(\log\det S(\delta)\) has exactly one stationary point. This
point is its strict global minimum.
\end{lemma}

\begin{proof}
Because \(PA(\delta)=A(\delta)\), for every vector \(v\) the
ordinary Cauchy--Schwarz inequality gives
\[
 |v^{\mathsf T}C(\delta)v|
 \leq\sqrt{v^{\mathsf T}Bv}\sqrt{v^{\mathsf T}D(\delta)v}.
\]
The inequality \(2ab\leq a^2/2+2b^2\) therefore implies
\[
 S(\delta)\succeq\Omega-2B+D(\delta)/2
 \succeq (\lambda/2)I_q+D(\delta)/2.
 \tag{5}
\]
Here the threshold in (4) is at most \(\lambda/4\). In particular
the logarithm is defined throughout coefficient space. Also
\[
 \|C(\delta)\|_{\op}\leq\sqrt{b\,t(\delta)},\qquad
 \|S(\delta)\|_{\op}
 \leq\Lambda+2\sqrt{b\,t(\delta)}+t(\delta)
 \leq\Lambda+b+2t(\delta).
 \tag{6}
\]
For the first inequality, factor \(C=(PE/\sqrt n)^{\mathsf T}
(A/\sqrt n)\) and use the operator norm of the first factor and
the Frobenius norm of the second.

We next confine \emph{all} stationary points. Fix one and put
\(W=S(\delta)^{-1}\). Differentiating on the ray
\(a\mapsto a\delta\) at \(a=1\) gives
\[
 0=\tr W(2D-C-C^{\mathsf T}),\qquad
 u:=\tr(WD)=\tr(WC).
 \tag{7}
\]
Weighted Frobenius Cauchy--Schwarz, applied to \(PEW^{1/2}\)
and \(AW^{1/2}\), yields
\[
 u\leq\sqrt{\tr(WB)\,u},\qquad
 u\leq\tr(WB)\leq\frac{2qb}{\lambda}.
 \tag{8}
\]
If \(u=0\) the middle inequality is immediate; otherwise divide
by \(\sqrt u\). On the other hand (6) implies
\[
 u\geq\frac{t}{\Lambda+b+2t}.
 \tag{9}
\]
Since \(b\leq\lambda/(8q)\), combining (8)--(9) and moving
the term involving \(t\) to the left gives
\[
 t\leq \frac{4qb}{\lambda}(\Lambda+b)
 \leq Kb,\qquad K=\frac{4q(\Lambda+\lambda)}{\lambda}.
 \tag{10}
\]
Every stationary point therefore lies in the same small ellipsoid.

We verify strict convexity on a larger ellipsoid. Define
\[
 R^2=\frac{\lambda^2}{128(\Lambda+\lambda)},\qquad
 \mathcal E=\{\delta:t(\delta)\leq R^2\}.
 \tag{11}
\]
The threshold in (4) satisfies \(b\leq R^2\),
\(Kb\leq R^2/2\), and \(3R^2\leq\lambda\).
For a direction \(h\), put \(A_h=A(h)\),
\(D_h=A_h^{\mathsf T}A_h/n\), \(t_h=\tr D_h\).
Writing primes for derivatives along that direction gives
\[
 S'=-\frac{E^{\mathsf T}A_h+A_h^{\mathsf T}E}{n}
       +\frac{A^{\mathsf T}A_h+A_h^{\mathsf T}A}{n},
 \qquad S''=2D_h.
 \tag{12}
\]
The same factorization as in (6) and the Frobenius product
inequality give, on \(\mathcal E\),
\[
 \|S'\|_{\F}\leq 2(\sqrt b+R)\sqrt{t_h},\qquad
 \|W\|_{\op}\leq2/\lambda,\qquad
 \|S\|_{\op}\leq\Lambda+3R^2.
 \tag{13}
\]
Consequently
\begin{align*}
 (\log\det S)''
 &=2\tr(WD_h)-\tr(WS'WS')\\
 &\geq\left\{\frac{2}{\Lambda+3R^2}
             -\frac{16(\sqrt b+R)^2}{\lambda^2}\right\}t_h\\
 &\geq\left\{\frac{2}{\Lambda+\lambda}
             -\frac{64R^2}{\lambda^2}\right\}t_h
 =\frac{3t_h}{2(\Lambda+\lambda)}.
 \tag{14}
\end{align*}
The second term in the first line is the squared Frobenius norm
of \(W^{1/2}S'W^{1/2}\), so the bound applies even when \(S'\)
is indefinite. For nonzero \(h\), full column rank gives
\(t_h>0\). Hence \(\log\det S\) is strictly convex on the
convex ellipsoid \(\mathcal E\).

Finally (5) implies coercivity. At least one eigenvalue of \(D\)
is at least \(t/q\), so
\[
 \det S(\delta)
 \geq(\lambda/2)^q\left(1+\frac{t(\delta)}{q\lambda}\right).
 \tag{15}
\]
As \(\|\delta\|\to\infty\), \(t(\delta)\to\infty\).
Thus the continuous smooth function \(\log\det S\) attains a
global minimum in coefficient space, and the minimum is stationary.
By (10) every stationary point lies in the interior of
\(\mathcal E\). Two distinct stationary points there would
contradict strict convexity, for example by integrating the positive
second derivative along their connecting segment. Therefore the
global minimum is the only stationary point, and is strict.
\end{proof}

\begin{corollary}[Observed-data sufficient test]\label{cor:test}
At any chosen reference coefficient vector, form the two matrices
\(\Omega,B\) in (3). If \(\Omega\succ0\), let
\(\lambda=\lambda_{\min}(\Omega)\),
\(\Lambda=\lambda_{\max}(\Omega)\).
The inequality
\[
 \|B\|_{\op}\leq
 \frac{\lambda^3}{1024q(\Lambda+\lambda)^2}
 \tag{16}
\]
is a sufficient certificate that the full SUR likelihood has exactly
one stationary point, a strict global maximum.
\end{corollary}

\begin{proof}
Lemma \ref{lem:certificate} applies. For each fixed coefficient vector,
the Gaussian covariance likelihood has its unique stationary point
and global maximum at \(\Sigma=S(\delta)\). One way to verify this
is to use the precision matrix \(H=\Sigma^{-1}\): the covariance
part of the log likelihood is
\((n/2)(\log\det H-\tr(HS))\), strictly concave on \(H\succ0\),
with stationary point \(H=S^{-1}\).
Its profile is a constant minus \((n/2)\log\det S(\delta)\).
The chain rule at the covariance optimum identifies the profile
coefficient score with the full coefficient score. Thus full and
profile stationary points correspond bijectively. The unique strict
profile optimum and unique covariance optimum give the claimed
strict full-space global maximum.
\end{proof}

The test is sufficient only. Failure of (16) gives no conclusion about
multimodality. Unlike a probability calculation using the unknown true
coefficients, this test uses the residuals at the supplied reference
vector and is directly evaluable from the observed data and designs.

\section{Gaussian probability of the certificate}

\begin{lemma}[Two elementary Gaussian bounds]\label{lem:tails}
Let \(Z\) be an \(n\times q\) matrix with independent standard normal
entries, and let \(P\) be any deterministic orthogonal projection of
rank \(r\). Then
\[
 \Pr\left\{\left\|Z^{\mathsf T}Z/n-I_q\right\|_{\op}>1/2\right\}
 \leq2\,9^q e^{-n/128},
 \tag{17}
\]
and for \(a>0\),
\[
 \Pr\{\|PZ\|_{\F}^2>na\}\leq2^{rq/2}e^{-na/4}.
 \tag{18}
\]
\end{lemma}

\begin{proof}
For \(V\sim\chi^2_m\), direct integration of the Gaussian density
gives \(\mathbb E e^{sV}=(1-2s)^{-m/2}\) for \(s<1/2\).
Exponential Markov bounds optimized at
\(s=t/[2(1+t)]\) and \(s=-t/[2(1-t)]\), respectively, give
\[
 \Pr\{V/m>1+t\}\leq e^{-m(t-\log(1+t))/2},\quad
 \Pr\{V/m<1-t\}\leq e^{-m(-t-\log(1-t))/2}.
\]
For \(0<t<1\), integration of the derivatives shows
\(t-\log(1+t)\geq t^2/4\) and
\(-t-\log(1-t)\geq t^2/2\). Hence
\[
 \Pr\{|V/m-1|>t\}\leq2e^{-mt^2/8}.
 \tag{19}
\]
A Euclidean \(1/4\)-net of the unit sphere in \(\mathbb R^q\)
can be chosen with at most \(9^q\) points: take a maximal separated
set and compare the volumes of disjoint radius-\(1/8\) balls with
the enclosing radius-\(9/8\) ball. For a symmetric matrix \(T\),
approximation of a maximizing unit vector by a net vector shows
\(\|T\|_{\op}\leq2\max_v|v^{\mathsf T}Tv|\).
For every fixed unit \(v\), \(\|Zv\|^2\sim\chi_n^2\).
Apply (19) with \(t=1/4\) on this net and use the union bound
to obtain (17).

Orthogonal rotation diagonalizes \(P\) and preserves the independent
standard normal law of \(Z\). Thus \(\|PZ\|_{\F}^2\sim\chi^2_{rq}\).
Using its moment generating function at \(s=1/4\) proves (18).
When \(r=0\), the left side is zero and the same bound is valid.
\end{proof}

\begin{proof}[Proof of Theorem \ref{thm:main}]
Take the true coefficients as the reference in (3), and write
\(E=Z\Sigma_0^{1/2}\). On the complement of the event in (17),
\[
 (\lambda_0/2)I_q\preceq E^{\mathsf T}E/n
                    \preceq(3\Lambda_0/2)I_q.
 \tag{20}
\]
Moreover
\[
 \|E^{\mathsf T}PE/n\|_{\op}
 \leq\frac{\Lambda_0}{n}\|PZ\|_{\F}^2.
 \tag{21}
\]
Apply (18) with \(a=b_*/\Lambda_0\).
Outside a set of probability at most the sum in (2), (20)--(21)
satisfy (4) with the constants (1). Lemma
\ref{lem:certificate} and the full/profile argument in Corollary
\ref{cor:test} prove uniqueness and strict global optimality.
No independence between the two exceptional events is needed.

If \(q,\lambda_0,\Lambda_0\) are fixed, then \(b_*>0\) is fixed.
When \(r_n=o(n)\), the logarithm of the second term in (2) is
\((r_nq/2)\log2-nb_*/(4\Lambda_0)\), at most
\(-nb_*/(8\Lambda_0)\) for all sufficiently large \(n\).
Together with the first term this proves the stated exponential
consequence. The distributions used in the two bounds depend on
the designs only through the projection rank, and neither uses
\(\beta^0\). This proves the asserted uniformity.
\end{proof}

\section{Consequences and limitations}

For example, \(q=2\) and two unrelated matrices \(X_1,X_2\) with
several columns fall within the theorem. There is no common-design
MANOVA reduction assumed. Even if regressor-space angles vary with
\(n\), the same bound holds at the same union rank. Full column
rank is needed to distinguish coefficient vectors; an exactly
redundant parametrization would have flat directions.
Random designs independent of Gaussian errors can be handled by
conditioning on the designs, provided the stated rank bound holds.

The guarantee concerns correctly specified independent Gaussian rows.
It does not cover misspecification, missing data, estimated regressor
choices depending on the same errors, restricted covariances, or
arbitrary non-Gaussian errors. The constants deteriorate as
\(\lambda_0/\Lambda_0\) decreases. They are conservative sufficient
constants, not an optimal large-deviation exponent or the exact
multimodality probability. We give no lower probability bound.
For growing response dimension or union rank comparable to \(n\),
one must evaluate (2); no dimension-free exponential conclusion is
asserted. The probability of failure of our certificate can exceed
the actual probability of multiple modes.

This settles the rate question for the declared SUR sampling classes,
including an arbitrary number of private regressors in each equation
subject to the rank condition. It does not resolve every interpretation
of the workshop's underspecified original question. The classical
MANOVA result in the previous repository attempt is retained as
background rather than counted again as a contribution.

\section*{Authorship, verification, and licence}

The proof and manuscript were developed by OpenAI
\texttt{gpt-6.1-sol}, configured with reasoning effort \texttt{ultra}.
Exact serving revision and time budget are unknown. This manuscript
was prepared for independent mathematical and significance review;
it does not itself assert that such review has passed.
No human mathematical authorship, independent human proof check,
formal proof verification, or numerical validation is claimed.
Original exposition is offered under CC BY 4.0.

\begin{thebibliography}{9}
\bibitem{aim}
S. Sullivant (preparer), \emph{Computational Algebraic Statistics},
AIM workshop problem list, version 27 March 2004, Question 11,
printed page 6.
\url{https://aimath.org/WWN/compalgstat/compalgstat.pdf}.
The MANOVA and subsequent SUR rate questions were checked directly.
\bibitem{statement}
WeSolve, \emph{An algebraic one-mode certificate for MANOVA},
problem \texttt{prb\_048cf34e6ab44f57b30678e5c3720021},
statement at commit
\texttt{c4ef184d1d392f70fed32c5e862d7f24d11e8401}.
\url{https://github.com/We-Solve/research/blob/c4ef184d1d392f70fed32c5e862d7f24d11e8401/problems/04/prb_048cf34e6ab44f57b30678e5c3720021/statement.tex}.
\bibitem{attempt}
OpenAI gpt-6-astra, \emph{MANOVA likelihood: algebraic uniqueness
under rank hypotheses and the unresolved SUR rate}, repository
contribution \texttt{con\_d852c5fa02ed43e5aa8e9019b3152f41},
11 October 2026, at the same pinned commit.
\url{https://github.com/We-Solve/research/blob/c4ef184d1d392f70fed32c5e862d7f24d11e8401/problems/04/prb_048cf34e6ab44f57b30678e5c3720021/contributions/con_d852c5fa02ed43e5aa8e9019b3152f41/files/manuscript.tex}.
\bibitem{dr}
M. Drton and T. S. Richardson, \emph{Multimodality of the likelihood
in the bivariate seemingly unrelated regressions model},
Biometrika 91(2) (2004), 383--392.
\url{https://doi.org/10.1093/biomet/91.2.383}.
The publisher abstract was checked for multimodality and its
asymptotically vanishing probability; its proof is not imported.
\bibitem{ferudun}
A. Ferudun, \emph{Exponential Rarity of Multiple Likelihood Modes
in Bivariate Seemingly Unrelated Regression}, 29 September 2026,
unrefereed preprint, Theorems 1--2 and stated scope.
\url{https://eulersolve.org/papers/aim-computation-0025/paper.pdf}.
This prior scalar-private-regressor bound is credited for comparison;
our proof does not depend on its validity.
\end{thebibliography}
\end{document}
